\subsection{Maturity and Adoption}

Beyond the coarse divide between prototypes and commercial products, maturity appears uneven across workflow stages. Assistants embedded in IDEs and CI environments tend to exhibit stronger reliability guarantees, clearer telemetry, and better policy controls, while end-to-end “app generators” and agentic toolchains still face churn in model backends and dependency stability. Reported productivity gains are promising but heterogeneous: controlled and field studies find time savings and reduced cognitive load, yet effects depend on task complexity, developer expertise, and the need for architectural reasoning \cite{peng2023impact,becker2025measuring}. Adoption is strongly mediated by organizational governance. Interviews indicate enthusiasm tempered by concerns about data exposure, provenance, and compliance, which shape procurement choices such as on-premises deployment, audit logging, and approved prompt patterns \cite{klemmer2024using}. Security findings further motivate cautious rollout and targeted training, given evidence that AI-assisted code can introduce or preserve vulnerabilities if safeguards and review practices are weak \cite{perry2023users,asare2024user}. In sum, the most mature deployments couple tool capability with operational guardrails and continuous evaluation aligned with software engineering guidance syntheses \cite{hou2024large}.

\subsection{Target Roles and Usage Models}

A broader role spectrum is emerging as capabilities move upstream and downstream of coding. Early evidence from requirements engineering suggests that LLMs can aid in drafting, refinement, and traceability when paired with structured prompts and review checkpoints, thereby inviting participation from analysts and domain experts under human-in-the-loop regimes (Cheng, 2024; generative). On the operations side, code-generating services for automation illustrate how DevOps teams adopt assistants through reproducible templates and policy-constrained suggestions that fit existing pipelines \cite{sahoo2024ansible}. Across roles, two usage models recur: co-creation, where practitioners iterate within familiar tools with immediate feedback and citations, and delegation, where agents orchestrate multi-step tasks under explicit constraints and post-hoc verification. Organizations that differentiate these modes, provide security-aware guidance, and instrument review workflows report smoother onboarding and more consistent outcomes \cite{klemmer2024using}.

% \subsection{Challenges and Limitations}
% Significant limitations include context-window restrictions, hallucinations, incomplete dependency resolution, and the requirement of post-editing by users.


% Ethical and legal concerns persist around code provenance, data privacy, and intellectual property.
